\documentclass[11pt,twocolumn]{article}
\usepackage[utf8]{inputenc}
\usepackage{amsmath,amssymb,mathtools,physics}
\usepackage{graphicx}
\usepackage{booktabs,tabularx,longtable}
\usepackage{tikz}
\usepackage{tcolorbox}
\usepackage{algorithm2e}
\usepackage{siunitx}
\usepackage[numbers]{natbib}
\usepackage{hyperref}
\usepackage{xeCJK}
\setCJKmainfont{SimSun}

\newcommand{\R}{\mathbb{R}}
\newcommand{\norm}[1]{\left\lVert #1 \right\rVert}
\DeclareMathOperator*{\argmin}{arg\,min}
\newenvironment{note}{\begin{tcolorbox}[colback=gray!5]}{\end{tcolorbox}}

\begin{document}
\title{复杂排版论文样例}
\author{Test}
\maketitle

\section{引言}
本文讨论范数 $\norm{x}_2$ 与集合 $x \in \R^n$ 的关系，参考 \cite{ref1}。
行内还有一个 \(a^2+b^2=c^2\) 的写法。

\subsection{方法}

这里给出一个行间公式：
\[
  \mathcal{L}(\theta) = \sum_{i=1}^{n} \left( y_i - f(x_i;\theta) \right)^2
\]

以及一段带 \verb|\verb| 与 \texttt{等宽} 的内容。

\subsubsection{子节}
再来一个带括号编号的公式：
\begin{equation}
  J = \int_{0}^{1} \norm{\nabla u}^2 \, dx
\end{equation}

双美元符号写法：
$$ p(x) = \frac{1}{\sqrt{2\pi\sigma^2}} e^{-\frac{(x-\mu)^2}{2\sigma^2}} $$


\begin{equation}
  \label{eq:main}
  \argmin_{\theta} \frac{1}{2}\norm{y - X\theta}_2^2 + \lambda \norm{\theta}_1
  \tag{1}
\end{equation}

\begin{align}
  a &= b + c \label{eq:a}\\
  d &= e - f \label{eq:b}
\end{align}

\begin{note}
这是一个 tcolorbox 装饰框。
\end{note}

\begin{figure}[h]
  \centering
  \begin{tikzpicture}
    \draw (0,0) -- (1,1);
  \end{tikzpicture}
  \caption{一个 TikZ 图}
\end{figure}

\begin{table}[h]
  \centering
  \begin{tabular}{lcc}
    \toprule
    参数 & 取值 & 单位 \\
    \midrule
    $\alpha$ & \SI{9.81}{\meter\per\second\squared} & -- \\
    \bottomrule
  \end{tabular}
  \caption{参数表}
\end{table}

\begin{algorithm2e}
  \KwIn{数据 $X$}
  \KwOut{参数 $\theta$}
  初始化 $\theta_0$ \;
  \For{$t=1$ \KwTo $T$}{
    更新 $\theta_t$ \;
  }
\end{algorithm2e}

\bibliographystyle{plain}
\bibliography{refs}
\end{document}
